\documentclass[10pt,a4paper]{amsart}
\usepackage[margin=19mm]{geometry}
\usepackage{amsmath,amssymb,mathtools}
\usepackage[scr=rsfs,cal=euler]{mathalfa}
\usepackage{xfrac}
\usepackage[unicode,hidelinks,bookmarks=false]{hyperref}
\newcommand{\ie}{i.e.\ }
\newcommand{\cf}{cf.\ }
\newcommand{\resp}{respectively\ }
\newcommand{\Subsection}{\subsection}
\providecommand{\nobreakdash}{\nobreak\hbox{-}}
\setlength{\emergencystretch}{2em}
\multlinegap0pt
\newtheorem{comment}{Comment}
\usepackage{tikz}
\usepackage{amssymb}
\usepackage{enumerate}
\usepackage{float}
\newtheorem{lem}{Lemma}
\title{Intersections of the Ekedahl-Oort and Newton strata of $\mathcal{A}_{5}$}
\author{Steven R. Groen, Elvira Lupoian, Mychelle Parker}
\date{Source: arXiv:2509.19878v2 (CC BY 4.0)}
\begin{document}
\maketitle

\noindent\emph{Source and licence.} Intersections of the Ekedahl-Oort and Newton strata of $\mathcal{A}_{5}$. Version arXiv:2509.19878v2 (CC BY 4.0), distributed by arXiv under the Creative Commons Attribution 4.0 International licence, http://creativecommons.org/licenses/by/4.0/, as recorded in the arXiv licence metadata for this exact version and shown on the abstract page https://arxiv.org/abs/2509.19878v2. Full licence notice: \texttt{licenses/arxiv-2509.19878-CC-BY-4.0.txt}.\par\smallskip

\noindent\textit{Editorial scope.} Counted unit: the article's lem 2/5 (source0.tex lines 1151--1153). The article's complete original proof of it is delivered with it (source0.tex lines 1154--1228, one \texttt{proof} environment of 75 lines). No mathematics is written, repaired or re-derived: the statement and the proof are the article's own lines, in the article's own notation, and no retained line is substituted or re-flowed.  Retained uncounted as the source-local context the proof needs: lines 1023--1025.  Nothing was omitted from any retained span, so the package carries no omission record.  No retained line contains a citation, so the wrapper carries no bibliography and no bibliography entry.  No retained line contains a figure environment, an image-inclusion command or drawing code, so no image is delivered and the package declares no asset dependency.  Editorial formatting only, in addition: an AMS \texttt{amsart} preamble that restates the article's own declarations and macros used by the retained lines, namely the environments \texttt{lem} on separate counters, together with the standard packages those lines need.  The retained environments print in this document as Lemma 1 in order of appearance, whereas the article prints the same items with the numbering of its own full document.  The complete native source of the article as distributed in the arXiv e-print for this exact version is bundled unchanged as \texttt{sources/185639/source0.tex}.
\begin{lem} \label{lem: Zink+} 
For every $n\in \mathbb{Z}_{\geq 1}$, we have $s_M^{(n)} = \frac{1}{n}\min_{i,j} \left\{v_p\left(f_{i,j}^{(n)}\right)\right\}$ and $s_M \geq s_M^{(n)}$.
\end{lem}

\begin{lem} \label{lem: (0,1,1,2,2) with 2/5}
The intersection $S_{(0,1,1,2,2)} \cap \mathcal{N}_{2/5}$ is non-empty.
\end{lem}

\begin{proof}
We now define the $\sigma$-linear operator $F$ on $M^\circ=W(k)\{a_1,\ldots a_{10} \}$ as follows:
\begin{comment}
\begin{equation*}
\begin{split}
\begin{aligned}
    F(a_1)&=a_2 \\
    F(a_2)&=a_3 \\
    F(a_3)&=a_4 \\
    F(a_4)&=pa_5 \\
    F(a_5)&=a_6+pa_1
\end{aligned}
\end{split}
\quad \quad
\begin{split}
\begin{aligned}
    F(a_6)&=pa_7 \\
    F(a_7)&=pa_8 \\
    F(a_8)&=pa_9 \\
    F(a_9)&=a_{10} \\
    F(a_{10})&=-pa_1
\end{aligned}
\end{split}
\end{equation*}
\end{comment}
\begin{equation*}
\begin{split}
\begin{aligned}
    F(a_1)&=a_2 \\
    F(a_6)&=pa_7
\end{aligned}
\end{split}
\quad \quad \quad
\begin{split}
\begin{aligned}
    F(a_2)&=a_3     \\
    F(a_7)&=pa_8
\end{aligned}
\end{split}
\quad \quad \quad
\begin{split}
\begin{aligned}
    F(a_3)&=a_4    \\
    F(a_8)&=pa_9 
\end{aligned}
\end{split}
\quad \quad \quad
\begin{split}
\begin{aligned}
    F(a_4)&=pa_5     \\
    F(a_9)&=a_{10}
\end{aligned}
\end{split}
\quad \quad \quad
\begin{split}
\begin{aligned}
    F(a_5)&=a_6+pa_1     \\
    F(a_{10})&=-pa_1.
\end{aligned}
\end{split}
\end{equation*}

As before, we let $V=pF^{-1}$ and set $\langle a_i,a_j\rangle=\delta_{i+5,j}$ for $1\leq i \leq 5$. We verify $A_F^TH=HA_V$ and conclude we have constructed a principally quasi-polarised free Dieudonn\'{e} module $M$. By reducing $A_F$ and $A_V$ modulo $p$, we establish $M/pM=\overline{M}(fffvfvvvfv) \cong N_{(0,1,1,2,2)}.$ Finally, we compute the Newton polygon of $M$. An explicit computation gives the block matrix
\[
A_F^{(5)} = \begin{bmatrix}
    \mathrm{diag}(p^2, p^2, p^2, p^2, p^2) & \mathrm{diag}(-p^4, -p^3, -p^2, -p, -p^2) \\
    \mathrm{diag}(p, p^2, p^3, p^4, p^3) & 0
\end{bmatrix}.
\]
We verify that $p^7 \mid A_F^{(20)}$ and hence $s_M\geq 7/20$ by Lemma~\ref{lem: Zink+}. On the other hand, inductively we show that  $v_p(f_{1,1}^{(5n)})=2n$ for every $n>0$, which implies $s_M \leq 2/5$. Hence we conclude $s_M=2/5$.
\begin{comment}
It is then straightforward to prove, using induction, that, for every $n>0$, we have $\min\left\{ v_p\left(f_{i,j}^{(5n)}\right) \right\} = 2n-1$, giving the first slope $\displaystyle s_M= \lim_{n\to \infty} \frac{2n-1}{5n} = \frac{2}{5}$ by Lemma~\ref{lem: Zink+}. 
\end{comment}

\end{proof}

\end{document}
